% Mathématiques 5e — cours V3
% Repérage dans le plan et représentation de l'espace

\section{Objectifs}

\begin{itemize}
\item Lire et placer une abscisse sur une droite graduée.
\item Lire les coordonnées d'un point dans un repère orthogonal.
\item Placer un point de coordonnées données.
\item Mettre en relation plusieurs représentations en perspective cavalière.
\item Relier perspective cavalière et patron pour un pavé droit, un prisme droit ou un cylindre de révolution.
\item Calculer le volume d'un cube, d'un pavé droit et d'un prisme droit.
\item Convertir des unités usuelles de volume et de capacité.
\item Calculer l'aire d'un disque.
\item Calculer le volume d'un cylindre de révolution.
\end{itemize}

\section{1. Abscisse sur une droite graduée}

Une droite graduée possède :
\begin{itemize}
\item une origine ;
\item une unité de longueur ;
\item un sens de lecture.
\end{itemize}

L'abscisse d'un point est le nombre qui repère sa position.

Si :
\[
A(-2{,}5),
\]
le point $A$ est placé à l'abscisse :
\[
-2{,}5.
\]

\section{2. Lire une graduation}

Avant de lire une abscisse, il faut déterminer la valeur d'un intervalle.

Si l'écart entre :
\[
0
\]
et :
\[
1
\]
est partagé en cinq intervalles égaux, chaque intervalle vaut :
\[
\dfrac15=0{,}2.
\]

Un point situé trois graduations à droite de zéro a donc pour abscisse :
\[
0{,}6.
\]

\section{3. Repère orthogonal}

Un repère orthogonal du plan comporte deux axes perpendiculaires :
\begin{itemize}
\item l'axe horizontal des abscisses ;
\item l'axe vertical des ordonnées.
\end{itemize}

Leur point d'intersection est l'origine :
\[
O.
\]

\coursefigure{/images/mathematiques/5eme/reperage-plan.svg}{Repère orthogonal contenant plusieurs points dans différents quadrants.}{Les coordonnées d'un point se lisent toujours dans l'ordre : abscisse puis ordonnée.}

\section{4. Coordonnées d'un point}

Un point $A$ de coordonnées :
\[
(x;y)
\]
se note :
\[
A(x;y).
\]

Le premier nombre est l'abscisse.

Le second est l'ordonnée.

\begin{attention}
Les coordonnées :
\[
(2;-3)
\]
et :
\[
(-3;2)
\]
ne désignent pas le même point.
\end{attention}

\section{5. Lire les coordonnées}

Pour lire les coordonnées d'un point :
\begin{enumerate}
\item projeter le point verticalement sur l'axe des abscisses ;
\item lire l'abscisse ;
\item projeter horizontalement sur l'axe des ordonnées ;
\item lire l'ordonnée.
\end{enumerate}

Si le point possède :
\[
x=-2
\]
et :
\[
y=3,
\]
on écrit :
\[
\boxed{A(-2;3)}.
\]

\section{6. Placer un point}

Pour placer :
\[
B(4;-1{,}5),
\]
on se déplace d'abord jusqu'à l'abscisse :
\[
4,
\]
puis verticalement jusqu'à l'ordonnée :
\[
-1{,}5.
\]

\begin{methode}
Toujours traiter les deux coordonnées séparément avant de placer le point.
\end{methode}

\section{7. Quadrants}

Les axes découpent le plan en quatre régions.

Pour un point :
\[
(x;y),
\]
les signes de $x$ et $y$ indiquent sa position.

Par exemple :
\[
x>0,\quad y>0
\]
place le point en haut à droite.

Et :
\[
x<0,\quad y<0
\]
le place en bas à gauche.

\section{8. Représenter un solide}

Une feuille est plane ; un solide possède trois dimensions.

La perspective cavalière représente donc un objet de l'espace sur une surface plane.

Les arêtes parallèles dans le solide sont généralement représentées par des segments parallèles dans le dessin.

Les arêtes cachées peuvent être tracées en pointillés.

\begin{attention}
Une perspective n'est pas une photographie exacte et ne conserve pas toutes les longueurs ni tous les angles.
\end{attention}

\section{9. Solides au programme}

En 5e, on travaille notamment sur :
\begin{itemize}
\item le cube ;
\item le pavé droit ;
\item le prisme droit ;
\item le cylindre de révolution.
\end{itemize}

\coursefigure{/images/mathematiques/5eme/espace-solides.svg}{Pavé droit, prisme droit et cylindre représentés en perspective.}{Pour calculer un volume, il faut identifier la base et la hauteur adaptée au solide.}

\section{10. Patron d'un solide}

Un patron est une figure plane qui permet de reconstruire le solide par pliage.

Pour un pavé droit, le patron est constitué de six rectangles.

Pour un prisme droit, il comporte :
\begin{itemize}
\item deux bases identiques ;
\item des faces latérales rectangulaires.
\end{itemize}

Pour un cylindre, un patron comprend :
\begin{itemize}
\item deux disques ;
\item un rectangle correspondant à la surface latérale déroulée.
\end{itemize}

\section{11. Relier patron et perspective}

Pour reconnaître un patron, il faut identifier :
\begin{itemize}
\item les faces qui deviendront opposées ;
\item les faces qui deviendront voisines ;
\item les longueurs qui doivent coïncider lors du pliage.
\end{itemize}

\begin{methode}
On peut mentalement fixer une face puis imaginer le pliage des faces adjacentes autour d'elle.
\end{methode}

\section{12. Volume}

Le volume mesure l'espace occupé par un solide.

L'unité de référence du système métrique est :
\[
\text{m}^3.
\]

On utilise aussi :
\[
\text{dm}^3,\quad \text{cm}^3,\quad \text{mm}^3.
\]

\section{13. Volume d'un cube}

Pour un cube d'arête $a$ :
\[
\boxed{V=a^3}.
\]

\begin{exemple}
Pour :
\[
a=4\ \text{cm},
\]
on obtient :
\[
V=4^3=64\ \text{cm}^3.
\]
\end{exemple}

\section{14. Volume d'un pavé droit}

Pour un pavé de longueur $L$, largeur $\ell$ et hauteur $h$ :
\[
\boxed{V=L\ell h}.
\]

\begin{exemple}
Si :
\[
L=8\ \text{cm},\quad \ell=5\ \text{cm},\quad h=3\ \text{cm},
\]
alors :
\[
V=8\times5\times3=120\ \text{cm}^3.
\]
\end{exemple}

\section{15. Volume d'un prisme droit}

Le volume d'un prisme droit est :
\[
\boxed{V=A_{\text{base}}\times h}.
\]

La difficulté consiste donc souvent à calculer l'aire de la base.

\begin{exemple}
Un prisme a une base triangulaire d'aire :
\[
12\ \text{cm}^2
\]
et une hauteur :
\[
7\ \text{cm}.
\]

Son volume vaut :
\[
V=12\times7=84\ \text{cm}^3.
\]
\end{exemple}

\section{16. Aire d'un disque}

Pour un disque de rayon $r$ :
\[
\boxed{A=\pi r^2}.
\]

Pour :
\[
r=3\ \text{cm},
\]
on obtient :
\[
A=9\pi\ \text{cm}^2.
\]

Une valeur approchée est :
\[
A\approx28{,}3\ \text{cm}^2.
\]

\section{17. Volume d'un cylindre}

Un cylindre de révolution est un prisme au sens volumique : son volume se calcule comme aire de base fois hauteur.

Sa base est un disque.

Donc :
\[
V=\pi r^2h.
\]

\begin{propriete}
Pour un cylindre de rayon $r$ et de hauteur $h$ :
\[
\boxed{V=\pi r^2h}.
\]
\end{propriete}

\section{18. Exemple développé : cylindre}

Un cylindre possède :
\[
r=2\ \text{cm}
\]
et :
\[
h=10\ \text{cm}.
\]

L'aire de la base vaut :
\[
A_{\text{base}}=\pi\times2^2=4\pi\ \text{cm}^2.
\]

Le volume vaut :
\[
V=4\pi\times10=40\pi\ \text{cm}^3.
\]

Donc :
\[
\boxed{V=40\pi\ \text{cm}^3}
\]
et approximativement :
\[
\boxed{V\approx125{,}7\ \text{cm}^3}.
\]

\section{19. Volume et capacité}

Certaines égalités doivent être connues.

\[
1\ \text{dm}^3=1\ \text{L}.
\]

Donc :
\[
1\ \text{cm}^3=1\ \text{mL}.
\]

Et :
\[
1\ \text{m}^3=1000\ \text{L}.
\]

Ces relations relient mesure géométrique d'un volume et capacité d'un récipient.

\section{20. Conversion des volumes}

Entre deux unités de volume successives, le facteur est :
\[
1000.
\]

En effet :
\[
1\ \text{dm}=10\ \text{cm},
\]
donc :
\[
1\ \text{dm}^3
=
10^3\ \text{cm}^3
=
1000\ \text{cm}^3.
\]

\begin{attention}
Pour les volumes, on ne multiplie pas par $10$ entre deux unités successives, mais par :
\[
1000.
\]
\end{attention}

\section{21. Exemple de conversion}

Convertissons :
\[
2{,}5\ \text{dm}^3
\]
en :
\[
\text{cm}^3.
\]

On utilise :
\[
1\ \text{dm}^3=1000\ \text{cm}^3.
\]

Donc :
\[
2{,}5\times1000=2500.
\]

Ainsi :
\[
\boxed{2{,}5\ \text{dm}^3=2500\ \text{cm}^3}.
\]

\section{22. Exemple mêlant capacité et volume}

Un aquarium possède un volume intérieur de :
\[
48\ \text{dm}^3.
\]

Comme :
\[
1\ \text{dm}^3=1\ \text{L},
\]
sa capacité maximale est :
\[
\boxed{48\ \text{L}}.
\]

\section{23. Contrôler une formule}

Un volume doit être exprimé en unité cubique.

Un résultat comme :
\[
120\ \text{cm}^2
\]
pour un volume est donc incorrect.

De même, une aire de disque doit être exprimée en unité carrée :
\[
\text{cm}^2.
\]

\section{24. Méthode de repérage}

\begin{methode}
Dans un repère :
\begin{enumerate}
\item identifier l'origine ;
\item vérifier les graduations des deux axes ;
\item lire ou placer l'abscisse ;
\item lire ou placer l'ordonnée ;
\item écrire les coordonnées dans l'ordre $(x;y)$.
\end{enumerate}
\end{methode}

\section{25. Méthode pour un volume}

\begin{methode}
Pour calculer un volume :
\begin{enumerate}
\item identifier le solide ;
\item repérer la base et la hauteur utiles ;
\item choisir la formule ;
\item convertir les longueurs dans une unité commune ;
\item calculer ;
\item écrire une unité cubique ;
\item convertir en capacité si le contexte le demande.
\end{enumerate}
\end{methode}

\section{26. Erreurs fréquentes}

\begin{itemize}
\item Inverser abscisse et ordonnée.
\item Lire un repère sans vérifier la graduation.
\item Mesurer une longueur directement sur une perspective.
\item Confondre aire et volume.
\item Utiliser $2\pi r$ à la place de $\pi r^2$ pour l'aire d'un disque.
\item Oublier de mettre le rayon au carré dans le volume du cylindre.
\item Convertir des unités de volume comme des unités de longueur.
\item Confondre $\text{cm}^3$ et $\text{cm}^2$.
\end{itemize}

\section{À retenir}

\begin{aretenir}
Dans un repère :
\[
\boxed{A(x;y)}
\]
signifie d'abord abscisse, puis ordonnée.

Les volumes essentiels sont :
\[
\boxed{V_{\text{cube}}=a^3},
\]
\[
\boxed{V_{\text{pavé}}=L\ell h},
\]
\[
\boxed{V_{\text{prisme}}=A_{\text{base}}h},
\]
\[
\boxed{V_{\text{cylindre}}=\pi r^2h}.
\]

Pour le disque :
\[
\boxed{A=\pi r^2}.
\]

Enfin :
\[
\boxed{1\ \text{dm}^3=1\ \text{L}}
\]
et :
\[
\boxed{1\ \text{cm}^3=1\ \text{mL}}.
\]
\end{aretenir}
